The propagator inverts the vertex $\sprod{2}{2}$ on its image, since $\sprod{2}{2}$ is injective on the quotient $\Hn{2} / \ker \sprod{2}{2}$.
At level $L$ and fixed total momentum, $\sprod{2}{2}$ is a finite block of integral operators in the relative momentum.
Since this operator is compact, its inverse on the image is unbounded.
  By ``invertible on the quotient'' we mean injective with a dense image.
  The propagator then rests on three assumptions, which we state here and discuss below.
  First, at each level, $\sprod{2}{2}$ has no null vector at generic total momentum, the squared total momenta at which it has one forming a discrete set.
  At level $0$ this is a result for the vertex family of~\cite{Erbin:2026:SL2CVerticesClosed} at its sections with $\beta > \sqrt{2}$, with no real exceptional momentum at all (\cref{sec:mom:definite}).
  Above level $0$ it rests on a Fredholm property, which we assume, and on the exclusion on physical grounds of a null vector at every momentum, which a heuristic argument supports.
  Second, the sources on which the inverse acts, such as the quintic source of \cref{p:sec:hs:quintic}, lie in its image, and not only in the closure of its image.
  Third, the quotient and the propagator have a limit as the level grows, at generic total momentum.

The first assumption is not automatic for an analytic family of compact operators.
  Zero is never an isolated point of the spectrum of an infinite-rank compact operator when the operator is self-adjoint, as ours is at real momentum up to the fixed reflection involution of the relative momentum, and there are entire families of Hilbert--Schmidt operators of the same type with null vectors on a set of momenta accumulating at a finite point.
  We can recover the finite-dimensional intuition with a Fredholm form.
  We divide $\sprod{2}{2}$ at level $L$ on both sides by the square root of its level-$0$ block at a fixed reference momentum (of its absolute value on the functions odd in the relative momentum, which appear above level $0$ and on which we assume the block injective).
  We assume that the result is a bounded operator with a bounded inverse at real $P^2$ at level $0$, and a Fredholm operator of index zero depending analytically on $P^2$ near the real axis at every level.
  The analytic Fredholm theorem then leaves only two possibilities at each level: a null vector at every total momentum, or null vectors at a discrete set of values of $P^2$ only.
A proof that $\sprod{2}{2}$ has no null vector at a single total momentum would select the second case at that level.
  At level $0$ the definiteness stated above does it; above level $0$ the form is indefinite near zero momentum, as recalled below, so the definiteness cannot be used, and no other proof is known.
  In the first case the quartic vertex would be blind to a whole spacetime field of $2$-string states.
  The quotient removes these states from $\Sigma_2$ and the source $\vertex{4}(\Sigma_2, \Psi^2)$ does not reach them, so the quartic order is unaffected.
  The construction then requires the quintic and every higher vertex to vanish on them at every momentum (the kernel conditions of \cref{p:sec:hs:quintic,pc:sec:all:kernel}), imposed on a whole field rather than at isolated momenta.
  The composition $\sprod{2}{2}^{-1} \circ \sprod{3}{2}$ needs in addition the second assumption, in either case.
  We exclude this first case on physical grounds: a sector of $2$-string states to which the quartic vertex is blind at every momentum would have been seen before, and the following heuristic argument supports the exclusion.

The stubs of the vertices suppress the high-mass states on each leg by a factor which, being injective, creates no null vector at fixed level.
Lengthening the stubs adds to the vertex region a strip of surfaces on which a single-string state of every mass is exchanged, the sum over these masses converging at finite stub length.
  The added stub length acts as a Laplace variable that separates the intermediate masses.
  A null vector that remains a null vector as the stubs are lengthened must then decouple from the exchange of every mass level separately.
  No $2$-string state should decouple from the whole tower of single-string exchanges: at level $0$ we expect this from the mechanism that makes the level-$0$ form definite (\cref{sec:mom:definite}), and above level $0$ we assume it.
  Under this decoupling assumption, a kernel present at every momentum cannot be a conservation law of the vertex family.
  The argument shows this much, but not that a kernel at every momentum cannot be accidental, that is, present at every momentum without a conservation law behind it.
  This is the possibility we excluded on physical grounds above.
  In support, neither the mechanisms that guarantee zeros of string vertices (the dilaton theorem with ghost-number counting, BRST-exact strands, the universality of zero-momentum couplings) nor the worldsheet twist produce a null vector of $\sprod{2}{2}$ at the levels and momenta where the couplings are known~\cite{Erbin:2026:SL2CVerticesClosed}, although these mechanisms are not exhaustive.

The same argument constrains the real zeros, that is, the real exceptional momenta, if there are any.
  A change of stub length acts on $\sprod{2}{2}$ at fixed level as a Gaussian factor in the relative momentum on each side, with a rate proportional to the added length, together with the exchanges across the added strip.
  For a fixed operator, the Gaussian factor keeps the null vectors that decay faster than every Gaussian and loses those with Gaussian decay once the added length exceeds a threshold set by their rate.
  Under the decoupling assumption above, a real zero at some level is not a conserved sector of the vertex family: its null vector cannot stay fixed, in fixed local coordinates, while the stubs are lengthened.
  A real zero that moves along a curve in the momentum--stub plane with a rotating null vector is not excluded.
  This is the generic behavior of an analytic family in two real parameters.
  Above level $0$ the form is indefinite near zero momentum (at level $1$, for the vertex family of~\cite{Erbin:2026:SL2CVerticesClosed}, \cref{p:sec:solutions}), so no argument we have excludes isolated real zeros there.
  For this reason, we state the first assumption at generic momentum rather than at every momentum.

The second assumption is separate from the first.
  The kernel condition of the quintic order~\eqref{p:eq:hs5:kernel-condition} places the quintic source in the closure of the image of $\sprod{2}{2}$, which is all that the vanishing of the quintic vertex on the kernel gives.
  The assumption, which we have not established, is that the source lies in the image itself.
  Gaussian decay of the source in the relative momentum is not enough: already for a model Gaussian kernel, only a window of decay rates of a Gaussian source is admissible.
  The third assumption is the weak form of the statement that the kernel stabilizes as the level grows.
  A null vector at level $L$ survives at level $L + 1$ only if the new rows of the form annihilate it, and level $L + 1$ can produce null vectors of its own, so the kernels at successive levels are not nested.
  We do not assume that every null vector beyond some level is annihilated by all higher rows, only that the quotient and the propagator converge at generic momentum.
  No finite-level computation decides either statement.
